\documentclass[11pt,a4paper]{article}
\usepackage[utf8]{inputenc}
\usepackage[T1]{fontenc}
\usepackage[spanish]{babel}
\usepackage{geometry}
\usepackage{booktabs}
\usepackage{longtable}
\usepackage{xcolor}
\usepackage{hyperref}
\usepackage{bytefield}
\usepackage{fancyhdr}
\usepackage{tikz}
\usepackage{listings}
\usetikzlibrary{arrows.meta, positioning}

\geometry{margin=2.5cm}

\hypersetup{
    colorlinks=true,
    linkcolor=blue!60!black,
    urlcolor=blue!60!black,
}

\lstset{
    basicstyle=\ttfamily\small,
    breaklines=true,
    frame=single,
    backgroundcolor=\color{gray!10},
    xleftmargin=1em,
    framexleftmargin=0.5em,
}

\pagestyle{fancy}
\fancyhf{}
\fancyhead[L]{\small Protocolo UART}
\fancyhead[R]{\small Coheteros}
\fancyfoot[C]{\thepage}

\title{Protocolo de Comunicación UART\\[0.3em]
\large STM32H723 $\leftrightarrow$ ESP32}
\author{Coheteros}
\date{\today}

\begin{document}

\maketitle

Todos los datos son \textbf{little-endian}, estructuras empaquetadas, sin relleno. El header \texttt{0xCAFE} se transmite LSB primero: \texttt{0xFE}, \texttt{0xCA}. Footer siempre \texttt{0xBE}.

\begin{figure}[h]
\centering
\resizebox{\textwidth}{!}{%
\begin{tikzpicture}[
    node distance=4cm,
    box/.style={draw, rounded corners, thick, minimum width=2.5cm, minimum height=1.5cm, align=center, font=\sffamily\small\bfseries},
    smallbox/.style={draw, rounded corners, thick, minimum width=2cm, minimum height=1.2cm, align=center, font=\sffamily\small\bfseries},
    arr/.style={-{Stealth[length=3mm]}, thick},
    lbl/.style={font=\sffamily\scriptsize, align=center, fill=white, inner sep=1pt},
]

% Nodes
\node[box] (et) {Estación\\Terrestre};
\node[box, right=of et] (esp) {ESP32};
\node[box, right=of esp] (cv) {Controlador\\de Vuelo};
\node[smallbox, above=2cm of esp] (gps) {Módulo\\GPS};

% Telemetry: CV -> ESP32 -> ET (top arrows)
\draw[arr, blue!70!black] ([yshift=0.3cm]cv.west) -- node[lbl, above] {Telemetría 95\,B} node[lbl, below] {\tiny UART} ([yshift=0.3cm]esp.east);
\draw[arr, blue!70!black] ([yshift=0.3cm]esp.west) -- node[lbl, above] {Telemetría 95\,B} node[lbl, below] {\tiny LoRa} ([yshift=0.3cm]et.east);

% Commands: ET -> ESP32 -> CV (bottom arrows)
\draw[arr, red!70!black] ([yshift=-0.3cm]et.east) -- node[lbl, above] {\tiny LoRa} node[lbl, below] {Comandos 5\,B} ([yshift=-0.3cm]esp.west);
\draw[arr, red!70!black] ([yshift=-0.3cm]esp.east) -- node[lbl, above] {\tiny UART} node[lbl, below] {Comandos 5\,B} ([yshift=-0.3cm]cv.west);

% GPS: GPS -> ESP32 -> CV
\draw[arr, green!50!black] (gps.south) -- node[lbl, right, xshift=2pt] {NMEA} (esp.north);
\draw[arr, green!50!black, dashed] ([yshift=0.7cm]esp.east) -- node[lbl, above] {GPS 24\,B} node[lbl, below] {\tiny UART} ([yshift=0.7cm]cv.west);

\end{tikzpicture}%
}
\caption{Flujo de datos. Línea punteada = trama construida por ESP32.}
\end{figure}

% ============================================================================
\section{Telemetría (CV $\rightarrow$ ESP32 $\rightarrow$ ET)}

Paquete de \textbf{95 bytes}. El ESP32 lo reenvía por LoRa \textbf{sin modificación}. No tiene bytes CMD ni LEN -- solo sync word y footer.

\begin{longtable}{@{}ccllp{4.5cm}@{}}
\toprule
\textbf{Offset} & \textbf{Bytes} & \textbf{Tipo} & \textbf{Campo} & \textbf{Codificación} \\
\midrule
\endfirsthead
\toprule
\textbf{Offset} & \textbf{Bytes} & \textbf{Tipo} & \textbf{Campo} & \textbf{Codificación} \\
\midrule
\endhead
0  & 2 & \texttt{uint16} & Sync             & \texttt{0xCAFE} \\
2  & 4 & \texttt{uint32} & Tick             & ms (crudo) \\
6  & 4 & \texttt{int32}  & AccelX           & m/s\textsuperscript{2} $\times$ 100 \\
10 & 4 & \texttt{int32}  & AccelY           & m/s\textsuperscript{2} $\times$ 100 \\
14 & 4 & \texttt{int32}  & AccelZ           & m/s\textsuperscript{2} $\times$ 100 \\
18 & 4 & \texttt{int32}  & GyroX            & deg/s $\times$ 100 \\
22 & 4 & \texttt{int32}  & GyroY            & deg/s $\times$ 100 \\
26 & 4 & \texttt{int32}  & GyroZ            & deg/s $\times$ 100 \\
30 & 4 & \texttt{int32}  & MagX             & mGauss $\times$ 100 \\
34 & 4 & \texttt{int32}  & MagY             & mGauss $\times$ 100 \\
38 & 4 & \texttt{int32}  & MagZ             & mGauss $\times$ 100 \\
42 & 4 & \texttt{int32}  & PressurePa       & Pa $\times$ 100 \\
46 & 4 & \texttt{int32}  & TemperatureC     & \textdegree C $\times$ 100 \\
50 & 4 & \texttt{int32}  & Latitude         & deg $\times$ 10\textsuperscript{7} \\
54 & 4 & \texttt{int32}  & Longitude        & deg $\times$ 10\textsuperscript{7} \\
58 & 4 & \texttt{int32}  & GPSAltitude      & m $\times$ 100 \\
62 & 1 & \texttt{uint8}  & Satellites       & crudo \\
63 & 4 & \texttt{int32}  & BaroAltitude     & m $\times$ 100 \\
67 & 4 & \texttt{int32}  & BaroVelocity     & m/s $\times$ 100 \\
71 & 4 & \texttt{int32}  & VelX             & m/s $\times$ 100 \\
75 & 4 & \texttt{int32}  & VelY             & m/s $\times$ 100 \\
79 & 4 & \texttt{int32}  & VelZ             & m/s $\times$ 100 \\
83 & 4 & \texttt{uint32} & FaultFlags       & máscara de bits \\
87 & 4 & \texttt{int32}  & BatteryVoltage   & V $\times$ 100 \\
91 & 1 & \texttt{uint8}  & State            & enum \\
92 & 1 & \texttt{uint8}  & RelayState       & máscara de bits \\
93 & 1 & \texttt{uint8}  & LastCommand      & enum \\
94 & 1 & \texttt{uint8}  & SyncEnd          & \texttt{0xBE} \\
\bottomrule
\caption{Paquete de telemetría (95 bytes).}
\end{longtable}

\subsection{Pseudocódigo: recepción de telemetría en ESP32}

\begin{lstlisting}[language=C, caption={ESP32: recibir telemetría y reenviar por LoRa.}]
// Buffer circular UART. Buscar sync word 0xFE 0xCA.
// Cuando se detecta, leer 95 bytes totales (incluyendo sync).
// Verificar que byte[94] == 0xBE (footer).
// Si es valido, enviar los 95 bytes por LoRa sin modificar.
// Si el footer no coincide, descartar y seguir buscando sync.

#define TELEMETRY_SIZE 95

if (uart_byte == 0xFE) {
    buffer[0] = 0xFE;
    uart_read(buffer + 1, TELEMETRY_SIZE - 1);
    if (buffer[1] == 0xCA && buffer[94] == 0xBE) {
        lora_send(buffer, TELEMETRY_SIZE);
    }
}
\end{lstlisting}


% ============================================================================
\section{Comandos (ET $\rightarrow$ ESP32 $\rightarrow$ CV)}

El ESP32 reenvía los bytes crudos recibidos por LoRa directamente al UART. \textbf{Sin parseo, sin validación, sin reempaquetado.}

Trama de \textbf{5 bytes} (comandos del operador no llevan payload):

\begin{table}[h]
\centering
\begin{tabular}{@{}cccll@{}}
\toprule
\textbf{Offset} & \textbf{Bytes} & \textbf{Tipo} & \textbf{Valor} & \textbf{Descripción} \\
\midrule
0 & 1 & \texttt{uint8} & \texttt{0xFE} & Sync LSB \\
1 & 1 & \texttt{uint8} & \texttt{0xCA} & Sync MSB \\
2 & 1 & \texttt{uint8} & CMD            & Tipo de comando \\
3 & 1 & \texttt{uint8} & \texttt{0x00}  & Longitud payload (cero) \\
4 & 1 & \texttt{uint8} & \texttt{0xBE}  & Footer \\
\bottomrule
\end{tabular}
\caption{Trama de comando (5 bytes).}
\end{table}

Valores válidos de CMD para comandos del operador:

\begin{table}[h]
\centering
\begin{tabular}{@{}cl@{}}
\toprule
\textbf{CMD} & \textbf{Acción} \\
\midrule
\texttt{0x01} & Reset \\
\texttt{0x02} & Aborto en tierra \\
\texttt{0x03} & Calibración \\
\texttt{0x04} & Drogue \\
\texttt{0x05} & Aterrizado \\
\bottomrule
\end{tabular}
\caption{Comandos del operador.}
\end{table}


% ============================================================================
\section{Datos GPS (GPS $\rightarrow$ ESP32 $\rightarrow$ CV)}

Única trama que el ESP32 \textbf{construye}. Parsea la sentencia NMEA \textbf{GGA} (\texttt{\$GPGGA} o \texttt{\$GNGGA}) del módulo GPS y empaqueta en una trama de comando con tipo \texttt{0x20}. \textbf{24 bytes} totales.

\begin{table}[h]
\centering
\begin{tabular}{@{}cccll@{}}
\toprule
\textbf{Offset} & \textbf{Bytes} & \textbf{Tipo} & \textbf{Valor} & \textbf{Descripción} \\
\midrule
0  & 1 & \texttt{uint8} & \texttt{0xFE}  & Sync LSB \\
1  & 1 & \texttt{uint8} & \texttt{0xCA}  & Sync MSB \\
2  & 1 & \texttt{uint8} & \texttt{0x20}  & Comando GPS \\
3  & 1 & \texttt{uint8} & \texttt{0x13}  & Longitud payload (19) \\
4--22 & 19 & --         & --            & Payload GPS \\
23 & 1 & \texttt{uint8} & \texttt{0xBE}  & Footer \\
\bottomrule
\end{tabular}
\caption{Trama GPS (24 bytes).}
\end{table}

\subsection{Payload GPS (19 bytes)}

\begin{table}[h]
\centering
\begin{tabular}{@{}ccclll@{}}
\toprule
\textbf{Offset} & \textbf{Bytes} & \textbf{Tipo} & \textbf{Campo} & \textbf{Unidad} & \textbf{Notas} \\
\midrule
0  & 4 & \texttt{uint32} & UnixTime     & Segundos      & Segundos del día (de GGA) \\
4  & 2 & \texttt{uint16} & Milliseconds & ms            & 0--999 \\
6  & 4 & \texttt{int32}  & Latitude     & deg $\times 10^7$ & \\
10 & 4 & \texttt{int32}  & Longitude    & deg $\times 10^7$ & \\
14 & 4 & \texttt{int32}  & Altitude     & mm            & Sobre nivel del mar \\
18 & 1 & \texttt{uint8}  & Satellites   & Cantidad      & \\
\bottomrule
\end{tabular}
\caption{Payload GPS (19 bytes).}
\end{table}

\subsection{Pseudocódigo: construcción de trama GPS en ESP32}

\begin{lstlisting}[language=C, caption={ESP32: parsear GGA y construir trama GPS.}]
// Sentencia GGA:
// $GNGGA,134523.500,4042.7688,N,07400.3600,W,1,08,1.0,30.5,M,-34.0,M,,*5E
//        ^^^^^^^^^^  ^^^^^^^^^^ ^ ^^^^^^^^^^^ ^ ^ ^^ ^^^ ^^^^
//        HHMMSS.sss  ddmm.mmmm   dddmm.mmmm     sat   alt(m)

// 1. Parsear hora: "134523.500"
uint32_t hours = 13, minutes = 45, seconds = 23, ms = 500;
uint32_t unix_time = hours * 3600 + minutes * 60 + seconds;
uint16_t milliseconds = ms;

// 2. Parsear latitud: "4042.7688,N"
//    dd = 40, mm.mmmm = 42.7688
//    grados = 40 + 42.7688 / 60 = 40.712813
//    N = positivo, S = negativo
int32_t latitude = (int32_t)(40.712813 * 1e7);  // 407128130

// 3. Parsear longitud: "07400.3600,W"
//    ddd = 074, mm.mmmm = 00.3600
//    grados = 74 + 00.3600 / 60 = 74.006000
//    W = negativo
int32_t longitude = (int32_t)(-74.006000 * 1e7);  // -740060000

// 4. Parsear altitud: "30.5" metros
int32_t altitude_mm = (int32_t)(30.5 * 1000);  // 30500

// 5. Parsear satelites: "08"
uint8_t satellites = 8;

// 6. Empaquetar payload (19 bytes, little-endian)
uint8_t payload[19];
memcpy(payload + 0,  &unix_time,    4);  // uint32 LE
memcpy(payload + 4,  &milliseconds, 2);  // uint16 LE
memcpy(payload + 6,  &latitude,     4);  // int32  LE
memcpy(payload + 10, &longitude,    4);  // int32  LE
memcpy(payload + 14, &altitude_mm,  4);  // int32  LE
payload[18] = satellites;

// 7. Construir trama completa (24 bytes)
uint8_t frame[24];
frame[0]  = 0xFE;          // sync LSB
frame[1]  = 0xCA;          // sync MSB
frame[2]  = 0x20;          // COMMAND_GPS_DATA
frame[3]  = 0x13;          // payload length = 19
memcpy(frame + 4, payload, 19);
frame[23] = 0xBE;          // footer

// 8. Enviar por UART al controlador de vuelo
uart_write(frame, 24);
\end{lstlisting}

% ============================================================================
\section{Resumen}

\begin{table}[h]
\centering
\begin{tabular}{@{}llcll@{}}
\toprule
\textbf{Datos} & \textbf{Dirección} & \textbf{Tamaño} & \textbf{ESP32} & \textbf{ET} \\
\midrule
Telemetría & CV $\rightarrow$ ET & 95\,B & Passthrough & Passthrough \\
Comandos   & ET $\rightarrow$ CV & 5\,B  & Passthrough & Passthrough \\
GPS        & GPS $\rightarrow$ CV & 24\,B & Construir trama & No involucrada \\
\bottomrule
\end{tabular}
\caption{CV = Controlador de Vuelo, ET = Estación Terrestre.}
\end{table}

\section{Estación Terrestre}

La estación terrestre es un passthrough bidireccional entre el módulo LoRa y una interfaz de usuario conectada por serial.

\begin{itemize}
    \item \textbf{Telemetría:} recibe los 95 bytes por LoRa y los escribe tal cual por serial para que la interfaz de usuario los lea y decodifique.
    \item \textbf{Comandos:} recibe los 5 bytes por serial desde la interfaz de usuario y los envía tal cual por LoRa.
\end{itemize}

La estación terrestre no parsea, valida ni modifica ningún paquete. Toda la lógica de decodificación y construcción de comandos está en la interfaz de usuario.

\end{document}
